%=======================================================================
%  INFORME — PRIMERA ENTREGA
%  Machine Learning · Universidad ESAN · Ciclo 2026-2
%  Seleccion Genomica de Precision en Ganado Vacuno
%
%  ARCHIVO UNICO. Se copia y se pega completo en Overleaf.
%  Compilador: pdfLaTeX. Compilar dos veces para que salgan las
%  referencias cruzadas (Overleaf lo hace solo).
%  No necesita ningun archivo externo: el diagrama de etapas esta
%  dibujado con TikZ y la bibliografia esta dentro del documento.
%
%  CONTENIDO EXIGIDO POR EL PROFESOR, EN ESTE ORDEN:
%    I.   Introduccion
%    II.  Estado del Arte  -> minimo 4 antecedentes, cada uno con
%         titulo/autores/anio, problema, BD, tecnicas, metodologia
%         y resultados. Aqui hay 5.
%    III. Metodologia propuesta -> grafico de etapas + descripcion
%         de cada etapa.
%    IV.  Experimentos -> la BD que se va a emplear.
%    Referencias
%
%  COMO ENCONTRAR LO QUE FALTA
%    Ctrl+F  >>> PENDIENTE   -> hay que escribirlo o verificarlo.
%    En el PDF sale en rojo entre corchetes.
%=======================================================================

\documentclass[conference]{IEEEtran}
\IEEEoverridecommandlockouts

\usepackage[T1]{fontenc}
\usepackage[utf8]{inputenc}
\usepackage[spanish,es-tabla,es-nodecimaldot,es-noquoting]{babel}
\usepackage{amsmath,amssymb}
\usepackage{booktabs}
\usepackage{array}
\usepackage{url}
\usepackage{xcolor}
\usepackage{tikz}
\usetikzlibrary{arrows.meta,positioning}
\usepackage[hidelinks]{hyperref}

% Colores del diagrama de etapas
\definecolor{cajaBorde}{HTML}{2F4858}
\definecolor{cajaFondo}{HTML}{EDF2F4}
\definecolor{cajaFondoAcc}{HTML}{DCE8E2}
\definecolor{textoTenue}{HTML}{5D6B66}

% Marca para texto que todavia falta escribir
\newcommand{\pendiente}[1]{\textcolor{red}{\textbf{[PENDIENTE: #1]}}}

\begin{document}

\title{Comparación de modelos de Machine Learning\\
frente a GBLUP para la predicción\\
del mérito genético en ganado vacuno\\
a partir de datos SNP de alta dimensionalidad}

% --- Bloque de autores: 3 arriba + 2 abajo ---------------------------
% Se arma con tablas y no con \and porque \and reparte los nombres
% segun el ancho disponible y los deja pegados unos con otros.
\author{%
\begin{tabular}{ccc}
Solange Lucero Ayala Mauricio & Luzmila Ana Cusi Apomayta & Joe José Mantilla Huamán \\
Facultad de Ingeniería        & Facultad de Ingeniería    & Facultad de Ingeniería \\
Universidad ESAN              & Universidad ESAN          & Universidad ESAN \\
Lima, Perú                    & Lima, Perú                & Lima, Perú \\
22200171@ue.edu.pe            & 22200284@ue.edu.pe        & 23100113@ue.edu.pe
\end{tabular}
\\[1.6em]
\begin{tabular}{cc}
Chayna Tucto Magiliano & Ivan Francisco Xavier Villegas Alarcón \\
Facultad de Ingeniería & Facultad de Ingeniería \\
Universidad ESAN       & Universidad ESAN \\
Lima, Perú             & Lima, Perú \\
% >>> PENDIENTE (IVAN): reemplazar XXXXXXXX por tu codigo de alumno.
23101292@ue.edu.pe     & XXXXXXXX@ue.edu.pe
\end{tabular}%
}

\maketitle

%=======================================================================
\begin{abstract}
La selección genómica permite estimar el mérito genético de un animal a
partir de su perfil de marcadores moleculares, sin esperar a que exprese
el rasgo ni a evaluar su descendencia. El método de referencia en la
industria es GBLUP, que asume efectos aditivos e infinitesimales sobre
todos los marcadores. Este trabajo evalúa si modelos de Machine Learning
capaces de representar efectos no aditivos superan a GBLUP en precisión
de predicción, y bajo qué condiciones lo hacen. Se trabaja con dos
fuentes genómicas de igual formato: el conjunto Holstein alemán de Zhang
\textit{et al.}, con 5\,024 toros genotipados en 42\,551 SNPs y tres
rasgos productivos, y un conjunto simulado con AlphaSimR en tres
escenarios de tamaño poblacional ($n = 125$, $500$ y $2\,000$) con
heredabilidad objetivo $h^2 = 0{,}30$. Los cinco modelos comparados son
GBLUP como línea base, regresión ridge, Random Forest, XGBoost y un
perceptrón multicapa, evaluados con particiones construidas sobre la
matriz de relación genómica para evitar fuga de información entre
parientes. La métrica principal es la correlación de Pearson entre el
valor predicho y el mérito genético observado. Los marcadores más
influyentes según SHAP se anotan a gen mediante la API REST de Ensembl y
se contrastan con la literatura mediante PubMed, de modo que la
interpretación del modelo quede respaldada por evidencia biológica
externa.
\end{abstract}

\begin{IEEEkeywords}
selección genómica; GBLUP; SNP; Machine Learning; XGBoost; Random Forest;
heredabilidad; validación cruzada genómica; SHAP; ganado vacuno.
\end{IEEEkeywords}

%=======================================================================
\section{Introducción}
%=======================================================================

El mejoramiento genético del ganado vacuno consiste en elegir, entre los
animales disponibles, aquellos cuya descendencia rendirá mejor. La
dificultad es que el rendimiento observable de un animal mezcla dos
componentes: lo que heredó y lo que le tocó de ambiente, alimentación,
manejo y clima. Solo el primero se transmite a la siguiente generación.
Durante décadas, separar ambos componentes exigió esperar a que el animal
produjera y a que su descendencia también lo hiciera, un ciclo que en
ganado lechero puede tomar cinco años o más por generación.

La selección genómica cambió ese esquema. Meuwissen, Hayes y Goddard
propusieron en 2001 estimar simultáneamente el efecto de decenas de miles
de segmentos cromosómicos a partir de un número limitado de registros
fenotípicos, y mostraron en simulación que con un mapa denso de
marcadores era posible predecir el valor genético de la descendencia con
una correlación de hasta 0{,}85, frente a 0{,}32 obtenido por mínimos
cuadrados \cite{meuwissen2001}. El mecanismo es el desequilibrio de
ligamiento: los marcadores cercanos a un gen de interés se heredan junto
con él, de modo que el marcador funciona como una etiqueta observable del
gen que no se observa. Un genotipo tomado al nacer pasa entonces a
contener información que antes exigía años de registros.

El estándar operativo derivado de esa idea es GBLUP, formulado en
términos computacionalmente tratables por VanRaden \cite{vanraden2008}.
GBLUP reemplaza el parentesco basado en pedigrí por una matriz de
relación genómica construida directamente a partir de los marcadores, y
resuelve un modelo mixto donde todos los marcadores contribuyen con
efectos pequeños provenientes de una misma distribución normal. Ese
supuesto, llamado infinitesimal, es a la vez la fortaleza y el límite del
método: lo vuelve estable y barato de calcular incluso con más marcadores
que animales, pero le impide representar dominancia, epistasis o
cualquier interacción entre marcadores.

Existen rasgos donde ese supuesto se sabe incorrecto. En ganado lechero,
el gen DGAT1 tiene un efecto documentado y grande sobre el contenido de
grasa de la leche, lo que contradice el supuesto infinitesimal para ese
rasgo en particular. Del mismo modo, cuando el rasgo depende de la
interacción entre el genotipo y el ambiente, un modelo puramente aditivo
sobre marcadores no dispone de la forma funcional necesaria para
representarla. Durbin \textit{et al.} documentaron en ganado vacuno
1\,040 asociaciones de interacción entre duración del día y genotipo
sobre la muda estacional de pelaje, evidencia directa de que la respuesta
de un mismo genotipo cambia según la señal ambiental \cite{durbin2024}.

Los modelos de Machine Learning sí pueden representar esas
interacciones. Sin embargo, la evidencia empírica en predicción genómica
no muestra una superioridad automática. Abdollahi-Arpanahi, Gianola y
Peñagaricano compararon seis métodos sobre 11\,790 toros Holstein
genotipados en 58\,000 SNPs y encontraron que gradient boosting alcanzó
la mejor correlación predictiva (0{,}36), seguido de Bayes B (0{,}34) y
GBLUP (0{,}33), mientras que las redes profundas quedaron por debajo
(0{,}29 para redes convolucionales y 0{,}26 para el perceptrón
multicapa); los métodos de aprendizaje profundo solo superaron a los
paramétricos cuando la varianza no aditiva era considerable y el tamaño
de muestra lo permitía \cite{abdollahi2020}. La diferencia entre el mejor
y el peor método fue de una décima de correlación, lo que indica que la
pregunta relevante no es cuál método es superior en abstracto, sino bajo
qué condiciones de tamaño poblacional y arquitectura genética conviene
uno u otro.

Ese es el vacío que aborda este trabajo. El objetivo general es comparar
el desempeño predictivo de modelos de Machine Learning frente al baseline
GBLUP en la estimación del mérito genético a partir de perfiles de
marcadores SNP, e identificar las condiciones de tamaño poblacional y
arquitectura genética bajo las cuales cada enfoque resulta preferible. La
pregunta de investigación que lo orienta es la siguiente: \textit{¿bajo
qué condiciones de tamaño poblacional, heredabilidad y arquitectura
genética los modelos de Machine Learning superan a GBLUP en la precisión
de predicción del mérito genético en ganado vacuno, cuando la validación
se construye respetando la estructura de parentesco de la población?}

Debe señalarse desde el inicio una restricción de los datos. En el Perú
no existe un conjunto genómico de ganado vacuno de acceso público con el
volumen que la selección genómica requiere, de modo que entrenar sobre
datos locales no es una opción disponible. La comparación se realiza por
lo tanto sobre dos fuentes complementarias: un conjunto real y público de
referencia internacional, que permite contrastar los resultados contra un
baseline ya publicado, y un conjunto simulado con parámetros controlados,
que permite variar deliberadamente el tamaño poblacional y la
heredabilidad para observar dónde cada modelo empieza a ganar o perder.
Las conclusiones del trabajo son, en consecuencia, metodológicas y no
constituyen recomendaciones de selección aplicables a un hato peruano.

%=======================================================================
\section{Estado del Arte}
%=======================================================================

Los cinco trabajos revisados se eligieron porque cada uno sostiene una
pieza concreta de este proyecto: el marco conceptual, el método base, el
conjunto de datos y su baseline publicado, la comparación contra modelos
de Machine Learning, y la evidencia de interacción genotipo por ambiente.
Cada uno se reporta con los seis campos requeridos.

%-----------------------------------------------------------------------
\subsection{Meuwissen, Hayes y Goddard (2001)}

\textbf{Título, autores y año.} \textit{Prediction of Total Genetic Value
Using Genome-Wide Dense Marker Maps}. Meuwissen, T.~H.~E., Hayes, B.~J.
y Goddard, M.~E. \textit{Genetics}, 157(4):1819--1829, 2001
\cite{meuwissen2001}.

\textbf{Problema.} Cómo estimar simultáneamente los efectos de decenas de
miles de haplotipos de marcadores a partir de un número limitado de
registros fenotípicos, cuando el número de efectos a estimar excede
ampliamente el número de observaciones disponibles.

\textbf{Base de datos.} Datos simulados. Un genoma de 1\,000~cM con
marcadores espaciados cada 1~cM, lo que produce aproximadamente 50\,000
haplotipos de marcadores, en una población de tamaño efectivo
$N_e = 100$. El desequilibrio de ligamiento entre marcadores y QTL surge
del tamaño finito de la población.

\textbf{Técnicas.} Mínimos cuadrados, BLUP de efectos de haplotipo con
varianza igual por segmento cromosómico, y los métodos bayesianos BayesA
y BayesB, que asumen una distribución previa sobre la varianza asociada a
cada segmento.

\textbf{Metodología.} Los marcadores que rodean cada región de 1~cM se
combinan en haplotipos. Sobre esos haplotipos se estiman los efectos con
cada método y se predice el valor genético de la descendencia de los
animales registrados. La comparación se realiza mediante la correlación
entre el valor genético predicho y el valor genético verdadero, conocido
por construcción al tratarse de simulación.

\textbf{Resultados.} Mínimos cuadrados alcanzó una precisión de 0{,}32,
por sobreestimación de los efectos incluidos. BLUP llegó a 0{,}73 pese a
que su supuesto de varianzas iguales era falso en la simulación. Los
métodos bayesianos elevaron la precisión a 0{,}85, incluso cuando la
distribución previa no correspondía a la verdadera.

%-----------------------------------------------------------------------
\subsection{VanRaden (2008)}

\textbf{Título, autores y año.} \textit{Efficient Methods to Compute
Genomic Predictions}. VanRaden, P.~M. \textit{Journal of Dairy Science},
91(11):4414--4423, 2008 \cite{vanraden2008}.

\textbf{Problema.} Cómo calcular predicciones genómicas de forma
computacionalmente eficiente cuando el número de marcadores excede el
número de animales, y cómo integrar la información de marcadores con la
del pedigrí tradicional.

\textbf{Base de datos.} Los algoritmos se derivaron y los programas se
probaron sobre datos simulados de 50\,000 marcadores y 2\,967 toros. La
validación sobre población real se publicó al año siguiente
\cite{vanraden2009}: 38\,416 marcadores tras control de calidad, con
3\,576 toros Holstein norteamericanos nacidos antes de 1999 como
conjunto predictor y 1\,759 toros nacidos entre 1999 y 2002 como
conjunto de predicción.

\textbf{Técnicas.} Tres formulaciones equivalentes de la predicción
genómica: una selección de marcadores, un modelo lineal sobre todos los
marcadores, y un modelo mixto sobre una matriz de relación genómica $G$.
La contribución central es la construcción de $G$ a partir de los
genotipos centrados.

\textbf{Metodología.} Los genotipos, codificados como número de copias
del alelo de referencia, se centran por la frecuencia alélica observada
para formar la matriz $Z$. La matriz de relación genómica se obtiene
entonces como
\begin{equation}
G = \frac{Z Z'}{2\sum_{j} p_j (1 - p_j)},
\label{eq:vanraden}
\end{equation}
donde $p_j$ es la frecuencia del alelo de referencia en el marcador $j$.
El denominador escala $G$ para que sea comparable con la matriz de
parentesco por pedigrí. Reemplazando el pedigrí por $G$ en las ecuaciones
del modelo mixto se obtiene GBLUP, cuyo costo depende del número de
animales y no del número de marcadores.

\textbf{Resultados.} Las tres formulaciones resultan equivalentes en las
predicciones que producen, pero difieren de manera sustancial en costo
computacional. El uso de $G$ permitió evaluaciones genómicas a escala
nacional y mejoró la confiabilidad de las predicciones respecto de las
basadas únicamente en pedigrí. La formulación se convirtió en el estándar
de la industria.

%-----------------------------------------------------------------------
\subsection{Zhang \textit{et al.} (2015)}

\textbf{Título, autores y año.} \textit{Accuracy of Whole-Genome
Prediction Using a Genetic Architecture-Enhanced Variance-Covariance
Matrix}. Zhang, Z., Erbe, M., He, J., Ober, U., Gao, N., Zhang, H.,
Simianer, H. y Li, J. \textit{G3: Genes, Genomes, Genetics},
5(4):615--627, 2015 \cite{zhang2015}.

\textbf{Problema.} GBLUP asume que todos los marcadores contribuyen con
efectos provenientes de una misma distribución. Cuando el rasgo tiene
marcadores de efecto grande, ese supuesto desperdicia información. El
problema es cómo incorporar conocimiento previo sobre la arquitectura
genética del rasgo dentro de la matriz de relación, sin abandonar la
eficiencia del modelo mixto.

\textbf{Base de datos.} 5\,024 toros Holstein alemanes genotipados en
42\,551 SNPs tras control de calidad, con valores fenotípicos
desregresados para tres rasgos: producción de leche en kilogramos,
porcentaje de grasa y score de células somáticas. Los datos están
disponibles como material suplementario (File~S1 y File~S2) del artículo.
El control de calidad —filtrado por frecuencia del alelo menor e
imputación de faltantes— viene ya aplicado por los autores.

\textbf{Técnicas.} GBLUP con la matriz $G$ estándar, BayesB, y el método
propuesto BLUP$|$GA, que construye una matriz de relación ponderada donde
un subconjunto de marcadores seleccionados recibe un peso mayor que el
resto.

\textbf{Metodología.} Validación cruzada con poblaciones de entrenamiento
de tres tamaños —2\,000, 500 y 125 animales— para cada uno de los tres
rasgos. El peso asignado al subconjunto de marcadores y el tamaño de ese
subconjunto se determinan mediante búsqueda sobre una rejilla de valores.
La precisión se mide como correlación entre el valor predicho y el
fenotipo desregresado del conjunto de prueba.

\textbf{Resultados.} BLUP$|$GA igualó o superó a GBLUP en los tres
rasgos, con ventaja mayor cuando el rasgo presentaba marcadores de efecto
grande y cuando la población de entrenamiento era pequeña. Las precisiones
por rasgo y por tamaño de entrenamiento están reportadas en el material
suplementario del artículo, lo que permite usarlas como referencia
externa.

\textbf{Relación con este proyecto.} Este es el conjunto de datos
principal del trabajo, y el artículo aporta algo más valioso que los
datos: una tabla de precisiones publicadas para los mismos tres tamaños
de entrenamiento. Esa tabla permite verificar que la implementación propia
de GBLUP es correcta antes de comparar cualquier modelo contra ella. Los
tres tamaños de población simulada de este proyecto —125, 500 y 2\,000—
se eligieron deliberadamente para coincidir con los del artículo.

%-----------------------------------------------------------------------
\subsection{Abdollahi-Arpanahi, Gianola y Peñagaricano (2020)}

\textbf{Título, autores y año.} \textit{Deep learning versus parametric
and ensemble methods for genomic prediction of complex phenotypes}.
Abdollahi-Arpanahi, R., Gianola, D. y Peñagaricano, F. \textit{Genetics
Selection Evolution}, 52:12, 2020 \cite{abdollahi2020}.

\textbf{Problema.} Si los métodos de aprendizaje profundo, capaces de
representar relaciones no lineales entre marcadores, superan a los
métodos paramétricos y de ensamble en predicción genómica de fenotipos
complejos, y bajo qué condiciones.

\textbf{Base de datos.} 11\,790 toros Holstein genotipados en 58\,000
SNPs, con la tasa de concepción del semental como rasgo objetivo.
Complementariamente, escenarios simulados construidos sobre los genotipos
reales observados, con heredabilidad de 0{,}30, efectos génicos aditivos
o no aditivos, y 100 o 1\,000 nucleótidos de rasgo cuantitativo.

\textbf{Técnicas.} Seis enfoques comparados sobre los mismos datos:
perceptrón multicapa, redes neuronales convolucionales, Random Forest,
gradient boosting, GBLUP y Bayes B.

\textbf{Metodología.} Validación cruzada sobre el conjunto real, y
repetición de la comparación sobre los escenarios simulados variando de
forma controlada la arquitectura genética —aditiva frente a no aditiva—
y el número de nucleótidos causales. La métrica de comparación es la
correlación predictiva entre valor predicho y valor observado.

\textbf{Resultados.} Sobre el conjunto real, el orden de correlaciones
predictivas fue: gradient boosting 0{,}36, Bayes~B 0{,}34, GBLUP 0{,}33,
Random Forest 0{,}32, redes convolucionales 0{,}29 y perceptrón multicapa
0{,}26. Los métodos de aprendizaje profundo superaron a los paramétricos
únicamente cuando la varianza no aditiva era considerable y el tamaño de
muestra resultaba suficiente.

\textbf{Relación con este proyecto.} Es el antecedente metodológico más
cercano y fija la expectativa realista del resultado: la diferencia entre
el mejor y el peor método fue de una décima de correlación, de modo que
un hallazgo honesto puede perfectamente consistir en que GBLUP no sea
superado. También orienta la selección de modelos: gradient boosting
—implementado aquí como XGBoost— es el candidato con mayor probabilidad
de superar al baseline, y el perceptrón multicapa se incluye como
contraste esperable a la baja en escenarios pequeños.

%-----------------------------------------------------------------------
\subsection{Durbin \textit{et al.} (2024)}

\textbf{Título, autores y año.} \textit{Genomic loci involved in sensing
environmental cues and metabolism affect seasonal coat shedding in Bos
taurus and Bos indicus cattle}. Durbin, H.~J., Yampara-Iquise, H., Rowan,
T.~N. \textit{et al.} \textit{G3: Genes, Genomes, Genetics},
14(2):jkad279, 2024 \cite{durbin2024}.

\textbf{Problema.} Qué loci del genoma bovino determinan la muda
estacional del pelaje, rasgo asociado a la tolerancia al calor, y en qué
medida el efecto de esos loci depende de señales ambientales como la
duración del día y la temperatura aparente.

\textbf{Base de datos.} 13\,364 animales con 36\,899 registros repetidos
de score de muda de pelaje, recolectados a lo largo de nueve años,
principalmente entre 2016 y 2019. Tras control de calidad e imputación se
retuvieron genotipos en 747\,009 marcadores para 11\,560 individuos. Los
datos son de acceso público en el repositorio Dryad.

\textbf{Técnicas.} Modelos mixtos con medidas repetidas para estimar
heredabilidad, estudios de asociación de genoma completo, y análisis de
interacción genotipo por ambiente sobre covariables ambientales
construidas.

\textbf{Metodología.} Para cada registro se calcularon la temperatura
aparente máxima media —que combina temperatura, humedad y viento— y la
duración media del día sobre los 30 días previos a la fecha de
evaluación. Esas covariables se incorporaron como término de interacción
con el genotipo, lo que permite distinguir un efecto génico constante de
uno que cambia según la señal ambiental.

\textbf{Resultados.} La heredabilidad en sentido estricto se estimó entre
0{,}321 y 0{,}407 según el subconjunto, con 0{,}368 (error estándar
0{,}012) en el conjunto completo. Se identificaron cientos de variantes
asociadas al rasgo, con señales fuertes en los cromosomas 5 y 23. El
análisis de interacción detectó 1\,040 asociaciones de duración del día
por genotipo y 17 de temperatura aparente por genotipo, lo que sitúa al
fotoperiodo como la señal ambiental dominante.

\textbf{Relación con este proyecto.} Aporta la evidencia empírica de que
el efecto de un marcador puede depender del ambiente, lo que sustenta la
hipótesis de que un modelo capaz de representar interacciones tiene margen
para superar a GBLUP. Además, su conjunto de datos aporta dos elementos
que el Holstein alemán no tiene: identificadores reales de SNP, que
permiten la anotación a gen, y covariables ambientales explícitas.

%-----------------------------------------------------------------------
\subsection{Síntesis}

La Tabla~\ref{tab:sintesis} resume los cinco trabajos. Dos patrones
orientan las decisiones de este proyecto. El primero es que GBLUP
constituye un baseline difícil de superar: se mantiene competitivo incluso
cuando sus supuestos son falsos, de modo que cualquier mejora debe
demostrarse sobre particiones que no favorezcan al modelo más flexible. El
segundo es que las ventajas de los modelos no lineales aparecen
condicionadas al tamaño de muestra y a la presencia efectiva de varianza
no aditiva, lo que convierte el tamaño poblacional en una variable
experimental y no en un dato del problema.

\begin{table*}[!t]
\centering
\caption{Síntesis de los antecedentes revisados.}
\label{tab:sintesis}
\renewcommand{\arraystretch}{1.25}
\footnotesize
\begin{tabular}{p{2.1cm} p{2.9cm} p{3.1cm} p{3.0cm} p{3.4cm}}
\toprule
\textbf{Trabajo} & \textbf{Problema} & \textbf{Base de datos} &
\textbf{Técnicas} & \textbf{Resultado principal} \\
\midrule
Meuwissen \textit{et al.} (2001) \cite{meuwissen2001} &
Estimar miles de efectos con pocos registros &
Simulación: 1\,000~cM, $\sim$50\,000 haplotipos, $N_e = 100$ &
Mínimos cuadrados, BLUP, BayesA, BayesB &
Precisión 0{,}32 / 0{,}73 / 0{,}85. Establece la viabilidad de la
selección genómica \\

VanRaden (2008) \cite{vanraden2008} &
Calcular predicciones genómicas de forma eficiente &
Simulación: 50\,000 marcadores, 2\,967 toros &
Tres formulaciones equivalentes; matriz $G$ &
Las tres son equivalentes en predicción pero difieren en costo; $G$ se
vuelve el estándar \\

Zhang \textit{et al.} (2015) \cite{zhang2015} &
Incorporar arquitectura genética en la matriz de relación &
5\,024 toros Holstein alemanes, 42\,551 SNPs, 3 rasgos &
GBLUP, BayesB, BLUP$|$GA &
BLUP$|$GA iguala o supera a GBLUP; mayor ventaja con entrenamiento
pequeño \\

Abdollahi-Arpanahi \textit{et al.} (2020) \cite{abdollahi2020} &
¿Superan las redes profundas a los métodos paramétricos? &
11\,790 toros Holstein, 58\,000 SNPs, más simulación con $h^2 = 0{,}30$ &
MLP, CNN, Random Forest, gradient boosting, GBLUP, Bayes~B &
Gradient boosting 0{,}36 $>$ Bayes~B 0{,}34 $>$ GBLUP 0{,}33 $>$ RF
0{,}32 $>$ CNN 0{,}29 $>$ MLP 0{,}26 \\

Durbin \textit{et al.} (2024) \cite{durbin2024} &
Loci de la muda estacional y su dependencia del ambiente &
13\,364 animales, 36\,899 registros, 747\,009 marcadores &
Modelos mixtos repetidos, GWAS, interacción G$\times$E &
$h^2 = 0{,}368$; 1\,040 interacciones duración del día $\times$ genotipo \\
\bottomrule
\end{tabular}
\end{table*}

%=======================================================================
\section{Metodología Propuesta}
%=======================================================================

El estudio es de tipo aplicado, cuantitativo y comparativo. Se organiza en
ocho etapas secuenciales, donde cada etapa entrega un artefacto
verificable que la siguiente consume. La Figura~\ref{fig:etapas} presenta
el flujo completo, indicando para cada etapa el artefacto que produce, y
las subsecciones que siguen describen cada una.

La separación en etapas con artefactos intermedios cumple dos funciones.
Permite que cinco personas trabajen en paralelo sin bloquearse entre sí, y
hace que un error se localice en la etapa que lo produjo en lugar de
aparecer al final como un resultado inexplicable.

% ---------------------------------------------------------------------
% DIAGRAMA DE ETAPAS
% Dibujado con TikZ, sin archivos externos. Si hay que mover una caja,
% cambiar su coordenada (x,y); las flechas siguen a los nodos solas.
% ---------------------------------------------------------------------
\begin{figure*}[!t]
\centering
% El resizebox ajusta el diagrama al ancho de la pagina. Si se agrega o
% se mueve una caja, no hay que recalcular nada: se reescala solo.
\resizebox{\textwidth}{!}{%
\begin{tikzpicture}[
  etapa/.style={
    draw=cajaBorde, line width=0.5pt, fill=cajaFondo, rounded corners=3pt,
    minimum width=3.5cm, minimum height=1.25cm, align=center,
    font=\footnotesize, inner sep=3pt},
  clave/.style={etapa, fill=cajaFondoAcc, line width=0.9pt},
  arte/.style={font=\scriptsize\itshape, text=textoTenue, align=center},
  arr/.style={-{Latex[length=2.2mm,width=1.6mm]}, draw=cajaBorde,
    line width=0.6pt}
]

% --- fila superior ---
\node[etapa]  (E1) at ( 0.0, 0) {\textbf{E1}\\Extracción y\\documentación de fuentes};
\node[etapa]  (E2) at ( 4.4, 0) {\textbf{E2}\\Preparación y\\contrato de carga};
\node[clave]  (E3) at ( 8.8, 0) {\textbf{E3}\\Particiones por\\relación genómica};
\node[etapa]  (E4) at (13.2, 0) {\textbf{E4}\\Baseline GBLUP\\y su validación};

\node[arte, below=1mm of E1] {6 fuentes documentadas};
\node[arte, below=1mm of E2] {\texttt{cargar() -> (X, y, ids)}};
\node[arte, below=1mm of E3] {\texttt{folds.csv}\\lo consumen E4 y E5 sin excepción};
\node[arte, below=1mm of E4] {precisión de referencia};

% --- fila inferior ---
\node[etapa]  (E5) at ( 0.0, -3.9) {\textbf{E5}\\Modelos de\\Machine Learning};
\node[etapa]  (E6) at ( 4.4, -3.9) {\textbf{E6}\\Evaluación\\comparativa};
\node[etapa]  (E7) at ( 8.8, -3.9) {\textbf{E7}\\Interpretabilidad y\\anotación biológica};
\node[etapa]  (E8) at (13.2, -3.9) {\textbf{E8}\\Dashboard\\interactivo};

\node[arte, below=1mm of E5] {4 modelos entrenados};
\node[arte, below=1mm of E6] {\texttt{resultados.csv}};
\node[arte, below=1mm of E7] {marcador $\to$ gen $\to$ evidencia};
\node[arte, below=1mm of E8] {visualización final};

% --- flechas de la fila superior ---
\draw[arr] (E1) -- (E2);
\draw[arr] (E2) -- (E3);
\draw[arr] (E3) -- (E4);

% --- retorno de E4 a E5: baja, cruza a la izquierda y entra por arriba ---
\draw[arr, rounded corners=5pt]
  (E4.south) -- (13.2,-2.15) -- (0,-2.15) -- (E5.north);

% --- flechas de la fila inferior ---
\draw[arr] (E5) -- (E6);
\draw[arr] (E6) -- (E7);
\draw[arr] (E7) -- (E8);

% --- E3 tambien alimenta a E5: las mismas particiones para todos ---
\draw[arr, dashed, draw=textoTenue]
  (8.8,-1.62) -- (8.8,-2.15);

\end{tikzpicture}%
}
\caption{Etapas de la metodología propuesta. Cada caja indica la etapa y,
debajo, el artefacto verificable que entrega. La etapa E3, resaltada,
genera el archivo de particiones que consumen tanto el baseline como los
cuatro modelos de Machine Learning, de modo que la comparación se realiza
siempre sobre las mismas divisiones de los datos.}
\label{fig:etapas}
\end{figure*}

%-----------------------------------------------------------------------
\subsection{Etapa E1 — Extracción y documentación de fuentes}

Se emplean seis fuentes agrupadas según la función que cumplen, detalladas
en la Tabla~\ref{tab:fuentes}. Las dos primeras comparten formato de
salida idéntico y son alternativas entre sí: entrenan el mismo modelo
sobre poblaciones distintas, pero no se fusionan, porque sus animales no
son los mismos. Las tres siguientes se encadenan aguas abajo del modelo,
no aguas arriba: reciben la lista de marcadores que el modelo identificó
como influyentes y devuelven evidencia biológica sobre ellos. La última
constituye un módulo separado, con su propia población y su propio
alcance.

Cada fuente entrega, además de los datos, un documento que permite a un
tercero repetir la descarga: enlace, licencia, número de filas y columnas,
separador, codificación y cualquier transformación aplicada. Los datos
crudos no se versionan por su tamaño; sí se versionan muestras reducidas
que permiten programar sin descargarlos.

\begin{table*}[!t]
\centering
\caption{Fuentes de datos del proyecto, su naturaleza y su función.}
\label{tab:fuentes}
\renewcommand{\arraystretch}{1.25}
\footnotesize
\begin{tabular}{p{2.5cm} p{2.0cm} p{1.5cm} p{4.3cm} p{4.5cm}}
\toprule
\textbf{Fuente} & \textbf{Tipo} & \textbf{API} & \textbf{Contenido} &
\textbf{Función en el proyecto} \\
\midrule
Holstein alemán \cite{zhang2015} & Estructurada & No &
5\,024 toros, 42\,551 SNPs, 3 rasgos productivos &
Conjunto de entrenamiento principal; su baseline publicado valida nuestra
implementación de GBLUP \\

Simulación AlphaSimR & Estructurada & No &
3 escenarios de 125, 500 y 2\,000 animales; 10\,000 SNPs; $h^2 = 0{,}30$ &
Control experimental: permite variar tamaño poblacional y heredabilidad de
forma deliberada \\

Mizzou Hair Shedding \cite{durbin2024} & Estructurada & No &
13\,364 animales, 36\,899 registros, 747\,009 marcadores, covariables
ambientales &
Referencia de interpretación: aporta identificadores reales de SNP y
evidencia de interacción G$\times$E \\

Ensembl REST & Semiestructurada (JSON) & Sí &
Posición cromosómica y gen más cercano por identificador de SNP &
Traduce los marcadores que el modelo destaca a nombres de gen \\

PubMed / E-utilities & No estructurada (texto) & Sí &
Resúmenes de artículos que mencionan cada gen &
Contrasta cada gen destacado con la literatura publicada \\

Cattle Side/Back View & No estructurada (imágenes) & No &
72 bovinos con fotografía lateral y posterior más medidas corporales &
Módulo complementario de visión, con población declarada independiente \\
\bottomrule
\end{tabular}
\end{table*}

%-----------------------------------------------------------------------
\subsection{Etapa E2 — Preparación y contrato de carga}

Toda fuente de entrenamiento se expone a través de la misma función,
independientemente de su origen:

\begin{center}
\texttt{cargar(...) -> (X, y, ids)}
\end{center}

donde \texttt{X} contiene los genotipos con \texttt{id\_animal} como
primera columna, \texttt{y} contiene los rasgos, e \texttt{ids} es el
vector de identificadores en el mismo orden que las filas de \texttt{X}.
La función valida antes de devolver: que la primera columna sea el
identificador, que los genotipos estén codificados en $\{0, 1, 2\}$, que
no haya valores faltantes ni identificadores duplicados, que el número de
filas coincida entre \texttt{X} e \texttt{y}, y que los identificadores de
ambas tablas sean idénticos en valor y en orden. Si alguna condición
falla, la carga se detiene con un mensaje que nombra el problema.

La última validación merece justificación. Un desalineamiento entre la
matriz de genotipos y la tabla de rasgos —un desplazamiento de una sola
fila— no produce ningún fallo visible: el modelo entrena sin errores y
devuelve correlaciones cercanas a cero, sin que nada indique el origen del
problema. Detectarlo en la carga es inmediato; detectarlo después de
comparar cinco modelos sobre tres escenarios obliga a repetir la etapa
completa.

Esta interfaz es el contrato que permite que los modelos sean agnósticos a
la fuente: el código de entrenamiento no distingue si los datos provienen
del conjunto alemán o de la simulación. Su cumplimiento está verificado
mediante pruebas automatizadas.

El control de calidad propiamente dicho depende de la fuente. Sobre el
conjunto Holstein no se repite, porque el filtrado por frecuencia del
alelo menor y la imputación vienen aplicados por los autores del artículo
original. Sobre el conjunto simulado la codificación se genera
directamente en $\{0, 1, 2\}$ sin faltantes. Sobre cualquier fuente que se
incorpore posteriormente se aplican dos filtros estándar: exclusión de
marcadores con frecuencia del alelo menor inferior a 0{,}01, por aportar
varianza despreciable, y exclusión de marcadores con más de 10\,\% de
datos faltantes; los faltantes remanentes se imputan por la media del
marcador.

%-----------------------------------------------------------------------
\subsection{Etapa E3 — Construcción de particiones de validación}

Esta es la decisión metodológica que más afecta la credibilidad de los
resultados. Los animales de una población de cría están emparentados entre
sí. Si la partición entre entrenamiento y prueba se realiza al azar, es
probable que un animal de prueba tenga un hermano completo en
entrenamiento, y el modelo obtendrá una correlación alta por reconocer al
pariente y no por haber aprendido la relación entre marcadores y rasgo. La
precisión reportada resulta entonces optimista y no representa el uso real
del método, que consiste en predecir sobre animales jóvenes de una
generación posterior.

El procedimiento adoptado es el siguiente. Primero se construye la matriz
de relación genómica $G$ según la Ecuación~\eqref{eq:vanraden}, usando
únicamente los genotipos y sin intervención de los rasgos. Luego se
agrupan los animales mediante agrupamiento jerárquico sobre la distancia
$1 - G_{ij}$, de modo que los animales con parentesco alto caigan en el
mismo grupo. Finalmente los grupos, no los animales individuales, se
reparten en cinco pliegues. Así se garantiza que ningún animal de prueba
comparta grupo con un animal de entrenamiento.

La asignación resultante se materializa en un único archivo,
\texttt{folds.csv}, con el identificador del animal y el número de pliegue
que le corresponde. Ese archivo se genera una sola vez y lo consumen todos
los modelos sin excepción, lo que garantiza que las diferencias observadas
entre modelos no sean diferencias entre particiones.

%-----------------------------------------------------------------------
\subsection{Etapa E4 — Baseline GBLUP y su validación}

Se implementa GBLUP como modelo mixto $y = \mu + g + e$ con
$g \sim N(0, G\sigma_g^2)$, donde $G$ es la matriz de la
Ecuación~\eqref{eq:vanraden}. Antes de compararlo con cualquier otro
modelo, su precisión se contrasta con los valores publicados por Zhang
\textit{et al.} \cite{zhang2015} para los mismos tres tamaños de población
de entrenamiento sobre el mismo conjunto de datos.

Esta verificación no es opcional. Si la implementación propia no reproduce
aproximadamente las cifras publicadas, cualquier diferencia posterior
entre GBLUP y un modelo de Machine Learning sería atribuible a un error de
implementación y no al modelo. El baseline se valida primero; la
comparación viene después.

%-----------------------------------------------------------------------
\subsection{Etapa E5 — Modelos de Machine Learning}

Se entrenan cuatro modelos sobre las particiones de la etapa E3:

\begin{itemize}
  \item \textbf{Regresión ridge.} Equivalente algebraico de GBLUP
    expresado sobre los marcadores en lugar de sobre los animales. Sirve
    de control interno: si ridge y GBLUP no coinciden, hay un error de
    implementación en alguno de los dos.
  \item \textbf{Random Forest.} Ensamble de árboles. Captura interacciones
    entre marcadores sin necesidad de especificarlas, pero fragmenta los
    efectos pequeños y distribuidos.
  \item \textbf{XGBoost.} Gradient boosting sobre árboles. Es el candidato
    con mayor probabilidad de superar al baseline según el antecedente de
    Abdollahi-Arpanahi \textit{et al.} \cite{abdollahi2020}.
  \item \textbf{Perceptrón multicapa.} Red neuronal densa. Se incluye para
    documentar el comportamiento del aprendizaje profundo en escenarios de
    muestra pequeña, donde la literatura anticipa desventaja.
\end{itemize}

Los hiperparámetros se ajustan mediante validación anidada dentro de cada
pliegue de entrenamiento, de modo que ningún dato del pliegue de prueba
participe en la selección de hiperparámetros.

%-----------------------------------------------------------------------
\subsection{Etapa E6 — Evaluación comparativa}

La métrica principal es la correlación de Pearson $r$ entre el valor
predicho y el valor de referencia del conjunto de prueba. Se usa
correlación y no $R^2$ porque en selección lo que importa es el orden
relativo de los animales: un modelo que ordena bien pero está desplazado
en escala sigue sirviendo para decidir a quién seleccionar. Como métrica
secundaria se reporta el RMSE, que sí es sensible a la escala. Se registra
además el tiempo de entrenamiento, porque un modelo que mejora la
correlación en 0{,}01 a costa de cien veces más cómputo puede no
justificar su adopción.

Cada combinación de modelo, rasgo, escenario y pliegue se escribe como una
fila en un único archivo de resultados con el esquema \texttt{modelo,
rasgo, escenario\_n, ruido, fold, r\_pearson, rmse, tiempo\_s}. Se reporta
la dispersión entre pliegues junto con el promedio, y no se declara
superioridad de un modelo sobre otro cuando los intervalos se solapan: con
aproximadamente mil animales de prueba por pliegue, diferencias de
correlación del orden de 0{,}01 pueden no ser distinguibles del error de
muestreo.

%-----------------------------------------------------------------------
\subsection{Etapa E7 — Interpretabilidad y anotación biológica}

Sobre el modelo de mejor desempeño se calculan valores SHAP para atribuir
la predicción a marcadores individuales. Los marcadores de mayor
atribución se consultan en la API REST de Ensembl para la especie
\textit{Bos taurus}, que devuelve cromosoma, posición y gen más cercano.
Cada gen obtenido se consulta luego en PubMed mediante la API E-utilities
del NCBI, contando menciones conjuntas del gen y de los términos del
rasgo.

El resultado es una tabla de tres columnas encadenadas —marcador, gen,
evidencia— que permite distinguir dos situaciones: un marcador destacado
por el modelo y respaldado por la literatura constituye evidencia
convergente; uno destacado sin respaldo queda señalado como hallazgo a
verificar. El gen DGAT1, con efecto documentado sobre el contenido de
grasa de la leche, se usa como caso de prueba del encadenamiento: si el
flujo funciona, ese gen debe aparecer con evidencia abundante.

Debe señalarse una restricción de la fuente principal. Los marcadores del
conjunto Holstein están etiquetados de forma secuencial, sin
correspondencia publicada con nombres reales de marcadores, de modo que la
anotación a gen no puede realizarse directamente sobre ellos. Por esa
razón el procedimiento se ilustra sobre la fuente que sí posee
identificadores reales.

%-----------------------------------------------------------------------
\subsection{Etapa E8 — Dashboard interactivo}

Los resultados se exponen en un dashboard que permite comparar modelos,
rasgos y escenarios, y consultar la evidencia biológica asociada a cada
marcador relevante. El dashboard lee el archivo de resultados producido en
la etapa E6 y no recalcula nada, de modo que lo que se muestra en pantalla
es exactamente lo que se midió.

El módulo complementario de visión se presenta en un apartado separado del
mismo dashboard. Ese módulo estima medidas corporales —peso, altura a la
cruz y perímetro torácico— a partir de fotografías lateral y posterior de
72 bovinos, con un pipeline independiente del genómico. El índice generado
contiene los 72 animales, todos con ambas vistas emparejadas por
identificador y sin valores faltantes, con cinco medidas corporales por
animal. Los 72 animales
fotografiados no tienen genotipos y los 5\,024 toros del conjunto genómico
no tienen fotografías: ambas poblaciones no comparten individuos y por lo
tanto no se fusionan en un mismo modelo. El módulo responde una pregunta
propia —si es posible caracterizar físicamente al animal sin balanza— y su
alcance se declara de forma explícita.

%-----------------------------------------------------------------------
\subsection{Herramientas}

El procesamiento de datos y el modelado se implementan en Python con
pandas, NumPy, scikit-learn y XGBoost \cite{xgboost}. La simulación de
poblaciones se realiza en R con el paquete AlphaSimR \cite{alphasimr}. La
interpretabilidad usa la biblioteca SHAP \cite{shap}. El control de
versiones es Git sobre GitHub, con una rama por integrante y revisión por
Pull Request antes de integrar a la rama principal. Las pruebas
automatizadas se ejecutan con pytest.

%=======================================================================
\section{Experimentos}
%=======================================================================

Esta sección presenta la base de datos que se emplea en el proyecto,
verificada sobre el contenido real de los archivos y no sobre la
descripción publicada.

%-----------------------------------------------------------------------
\subsection{Conjunto Holstein alemán}

Los archivos suplementarios File~S1 y File~S2 del artículo de Zhang
\textit{et al.} \cite{zhang2015} fueron descargados y verificados. El
archivo de genotipos contiene 5\,024 animales y 42\,551 marcadores, con
valores exclusivamente en $\{0, 1, 2\}$ y sin datos faltantes, ocupando
428~MB una vez descomprimido. El archivo de fenotipos contiene 5\,024
filas con las columnas de identificador y los tres rasgos, estandarizadas
a media cero y desviación unitaria, sin nulos. El material suplementario
incluye además una tabla con tres hojas —porcentaje de grasa con 1\,325
registros, producción de leche con 1\,622 y score de células somáticas con
933— y el documento con el baseline publicado para poblaciones de
entrenamiento de 2\,000, 500 y 125 animales.

La estructura de las dos tablas es la siguiente:

\begin{center}
\footnotesize
\begin{tabular}{@{}ll@{}}
\toprule
\textbf{Tabla} & \textbf{Columnas} \\
\midrule
Genotipos & \texttt{id\_animal}, \texttt{snp\_1} $\ldots$
  \texttt{snp\_42551} \\
Fenotipos & \texttt{id\_animal}, \texttt{mkg}, \texttt{fpro},
  \texttt{scs} \\
\bottomrule
\end{tabular}
\end{center}

donde \texttt{mkg} es producción de leche en kilogramos, \texttt{fpro} es
porcentaje de grasa y \texttt{scs} es el score de células somáticas. Los
genotipos se codifican como número de copias del alelo de referencia.

%-----------------------------------------------------------------------
\subsection{Conjunto simulado con AlphaSimR}

Se generaron tres escenarios con los parámetros de la
Tabla~\ref{tab:params-sim}, elegidos para que los tamaños coincidan con
los del baseline publicado y permitan la comparación directa. Los
resultados verificados sobre los archivos generados están en la
Tabla~\ref{tab:res-sim}.

\begin{table}[!t]
\centering
\caption{Parámetros de la simulación de poblaciones.}
\label{tab:params-sim}
\renewcommand{\arraystretch}{1.2}
\footnotesize
\begin{tabular}{@{}ll@{}}
\toprule
\textbf{Parámetro} & \textbf{Valor} \\
\midrule
Tamaños poblacionales & 125, 500 y 2\,000 animales \\
Marcadores por individuo & 10\,000 SNPs \\
Loci de rasgo cuantitativo (QTL) & 1\,000 \\
Heredabilidad objetivo & $h^2 = 0{,}30$ \\
Cromosomas & 29 (cariotipo bovino) \\
Especie base del coalescente & \texttt{CATTLE} \\
Arquitectura del rasgo & Puramente aditiva \\
Semilla & 22200285 \\
Herramientas & R 4.6.1, AlphaSimR 2.1.0 \\
\bottomrule
\end{tabular}
\end{table}

\begin{table}[!t]
\centering
\caption{Resultados de las tres corridas de simulación.}
\label{tab:res-sim}
\renewcommand{\arraystretch}{1.2}
\footnotesize
\begin{tabular}{@{}rrrr@{}}
\toprule
\textbf{$n$} & \textbf{$h^2$ realizada} & \textbf{Tiempo (s)} &
\textbf{Tamaño (MB)} \\
\midrule
125   & 0{,}283 & 148{,}4  & 2{,}6 \\
500   & 0{,}289 & 1\,098{,}6 & 9{,}7 \\
2\,000 & 0{,}304 & 3\,906{,}6 & 39{,}0 \\
\bottomrule
\end{tabular}
\end{table}

Dos observaciones resultan de estas corridas. La primera es que la
heredabilidad realizada converge al valor objetivo de 0{,}30 a medida que
crece $n$: con 125 animales se obtiene 0{,}283 y con 2\,000 se obtiene
0{,}304. La desviación en muestras pequeñas es error de muestreo, no un
defecto de la simulación, y constituye por sí misma una ilustración del
problema que el proyecto estudia: con pocos animales, incluso una cantidad
fijada por parámetro se estima con ruido. La segunda es que el costo
computacional crece de forma marcadamente no lineal, de dos minutos y
medio a más de una hora, lo que obliga a versionar la salida de la corrida
mayor en lugar de regenerarla en cada ejecución.

Cada escenario produce dos archivos con la estructura del contrato de
carga: la matriz de genotipos con \texttt{id\_animal} como primera
columna, y la tabla de rasgos con el fenotipo observado y el valor
genético verdadero. Este último solo existe por tratarse de simulación, y
es lo que permite medir la precisión de cada modelo contra la verdad
conocida y no contra un fenotipo ruidoso.

%-----------------------------------------------------------------------
\subsection{Verificación automatizada del contrato de datos}

Cada fuente se acompaña de pruebas que verifican el contrato antes de que
los datos lleguen al modelo. A la fecha hay 33 pruebas en ejecución, todas
en verde, que comprueban entre otras cosas: que los identificadores estén
alineados con las filas de la matriz de genotipos; que una tabla de rasgos
con las filas reordenadas sea rechazada; que un genotipo fuera del rango
$\{0, 1, 2\}$ sea rechazado; que un rasgo inexistente produzca un error
nombrado en lugar de un fallo silencioso; y que la ausencia de un
escenario simulado produzca un mensaje que indique el comando exacto para
generarlo.

%-----------------------------------------------------------------------
\subsection{Disponibilidad de los datos}

El conjunto de genotipos ocupa 428~MB, por lo que no se adjunta al archivo
de entrega ni se versiona en el repositorio. Se comparte mediante enlace,
junto con la documentación que permite repetir la descarga desde la fuente
original. En el repositorio se versionan muestras reducidas que permiten
ejecutar el código sin descargar el conjunto completo.

\begin{center}
\footnotesize
\begin{tabular}{@{}p{2.5cm}p{5.1cm}@{}}
\toprule
\textbf{Recurso} & \textbf{Ubicación} \\
\midrule
Repositorio del proyecto &
  \url{https://github.com/lulu1604/seleccion-genomica-ml} \\[2pt]
Conjunto de datos completo &
  \pendiente{ANA: pegar el enlace de Drive} \\[2pt]
Fuente original de los datos &
  \url{https://doi.org/10.1534/g3.114.016261} \\
\bottomrule
\end{tabular}
\end{center}

%=======================================================================
% REFERENCIAS
% Escritas a mano dentro del documento para que todo quepa en un solo
% archivo. Para agregar una: copiar un \bibitem, cambiar la clave y los
% datos, y citarla en el texto con \cite{clave}. Si no se cita, no
% aparece numerada correctamente.
%=======================================================================
\begin{thebibliography}{9}

\bibitem{meuwissen2001}
T.~H.~E. Meuwissen, B.~J. Hayes, and M.~E. Goddard,
``Prediction of total genetic value using genome-wide dense marker maps,''
\emph{Genetics}, vol.~157, no.~4, pp.~1819--1829, 2001.
doi: 10.1093/genetics/157.4.1819.

\bibitem{vanraden2008}
P.~M. VanRaden,
``Efficient methods to compute genomic predictions,''
\emph{Journal of Dairy Science}, vol.~91, no.~11, pp.~4414--4423, 2008.
doi: 10.3168/jds.2007-0980.

\bibitem{vanraden2009}
P.~M. VanRaden, C.~P. Van Tassell, G.~R. Wiggans, T.~S. Sonstegard,
R.~D. Schnabel, J.~F. Taylor, and F.~S. Schenkel,
``Invited review: Reliability of genomic predictions for North American
Holstein bulls,''
\emph{Journal of Dairy Science}, vol.~92, no.~1, pp.~16--24, 2009.
doi: 10.3168/jds.2008-1514.

\bibitem{zhang2015}
Z.~Zhang, M.~Erbe, J.~He, U.~Ober, N.~Gao, H.~Zhang, H.~Simianer, and
J.~Li,
``Accuracy of whole-genome prediction using a genetic
architecture-enhanced variance-covariance matrix,''
\emph{G3: Genes, Genomes, Genetics}, vol.~5, no.~4, pp.~615--627, 2015.
doi: 10.1534/g3.114.016261.

\bibitem{abdollahi2020}
R.~Abdollahi-Arpanahi, D.~Gianola, and F.~Peñagaricano,
``Deep learning versus parametric and ensemble methods for genomic
prediction of complex phenotypes,''
\emph{Genetics Selection Evolution}, vol.~52, art.~12, 2020.
doi: 10.1186/s12711-020-00531-z.

\bibitem{durbin2024}
H.~J. Durbin, H.~Yampara-Iquise, T.~N. Rowan \emph{et al.},
``Genomic loci involved in sensing environmental cues and metabolism
affect seasonal coat shedding in \emph{Bos taurus} and \emph{Bos indicus}
cattle,''
\emph{G3: Genes, Genomes, Genetics}, vol.~14, no.~2, art.~jkad279, 2024.
doi: 10.1093/g3journal/jkad279.

\bibitem{alphasimr}
R.~C. Gaynor, G.~Gorjanc, and J.~M. Hickey,
``AlphaSimR: an R package for breeding program simulations,''
\emph{G3: Genes, Genomes, Genetics}, vol.~11, no.~2, art.~jkaa017, 2021.
doi: 10.1093/g3journal/jkaa017.

\bibitem{shap}
S.~M. Lundberg and S.-I. Lee,
``A unified approach to interpreting model predictions,''
in \emph{Advances in Neural Information Processing Systems 30
(NIPS 2017)}, pp.~4765--4774, 2017.

\bibitem{xgboost}
T.~Chen and C.~Guestrin,
``XGBoost: A scalable tree boosting system,''
in \emph{Proceedings of the 22nd ACM SIGKDD International Conference on
Knowledge Discovery and Data Mining}, pp.~785--794, 2016.
doi: 10.1145/2939672.2939785.

\end{thebibliography}

\end{document}
